\chapter{Fundamentação Teórica}
\label{sec:fundamentacao}

Esta seção apresenta os conceitos fundamentais que sustentam este trabalho, incluindo a estrutura da certificação PMP, os princípios dos Modelos de Linguagem de Grande Porte (LLMs) e as estratégias de interação por meio de \textit{prompt engineering}. Ao final, são discutidos os desafios inerentes ao uso de LLMs em tarefas que envolvem julgamento contextual e tomada de decisão.




\section{Certificação PMP e Avaliação Baseada em Competências}
\label{subsec:pmp}

\textbf{Ideias a abordar:}
\begin{itemize}
    \item Definição da certificação PMP e seu papel na área de gerenciamento de projetos.
    \item Estrutura conceitual do PMP:
    \begin{itemize}
        \item Avaliação baseada em competências (não apenas conhecimento)
        \item Ênfase em aplicação prática
    \end{itemize}
    \item Organização do exame segundo o \textit{Exam Content Outline (ECO)}:
    \begin{itemize}
        \item Domínios: People, Process, Business Environment
    \end{itemize}
    \item Natureza das questões:
    \begin{itemize}
        \item Cenários situacionais
        \item Problemas abertos com múltiplas alternativas plausíveis
        \item Necessidade de julgamento e priorização
    \end{itemize}
    \item Integração de abordagens:
    \begin{itemize}
        \item Preditiva
        \item Ágil
        \item Híbrida
    \end{itemize}
    \item Implicação central:
    \begin{itemize}
        \item O exame exige raciocínio contextual e tomada de decisão
    \end{itemize}
\end{itemize}




\section{Modelos de Linguagem de Grande Porte (LLMs)}
\label{subsec:llms}

\textbf{Ideias a abordar:}
\begin{itemize}
    \item Definição de LLMs e base em arquiteturas do tipo Transformer.
    \item Princípio de funcionamento:
    \begin{itemize}
        \item Modelagem probabilística da linguagem
        \item Predição de tokens com base em contexto
    \end{itemize}
    \item Capacidades principais:
    \begin{itemize}
        \item Compreensão de texto
        \item Geração de respostas
        \item Raciocínio aproximado
    \end{itemize}
    \item Características relevantes:
    \begin{itemize}
        \item Sensibilidade ao contexto fornecido
        \item Dependência da formulação do prompt
    \end{itemize}
    \item Limitações fundamentais:
    \begin{itemize}
        \item Alucinações
        \item Falta de garantia de veracidade
        \item Ausência de modelo explícito de mundo
    \end{itemize}
\end{itemize}





\section{Prompt Engineering}
\label{subsec:prompt}

\textbf{Ideias a abordar:}
\begin{itemize}
    \item Definição de \textit{prompt engineering} como mecanismo de controle do comportamento do modelo.
    \item Papel do prompt como interface entre usuário e LLM.
    \item Principais estratégias:
    \begin{itemize}
        \item Zero-shot
        \item Few-shot
        \item Chain-of-Thought (CoT)
        \item Prompt com grounding (uso de contexto externo)
    \end{itemize}
    \item Impacto do prompt:
    \begin{itemize}
        \item Influência na qualidade da resposta
        \item Influência na estrutura do raciocínio
    \end{itemize}
    \item Limitação:
    \begin{itemize}
        \item Não elimina completamente erros ou inconsistências
    \end{itemize}
\end{itemize}





\section{Raciocínio e Tomada de Decisão em LLMs}
\label{subsec:decision}

\textbf{Ideias a abordar:}
\begin{itemize}
    \item Diferença entre:
    \begin{itemize}
        \item Recuperação de conhecimento (factual)
        \item Raciocínio (inferência)
        \item Tomada de decisão (escolha entre alternativas)
    \end{itemize}
    \item Natureza da tomada de decisão em gerenciamento de projetos:
    \begin{itemize}
        \item Dependência de contexto
        \item Trade-offs entre alternativas
        \item Ausência de resposta única trivial
    \end{itemize}
    \item Limitações dos LLMs nesse contexto:
    \begin{itemize}
        \item Tendência a respostas genéricas
        \item Dificuldade em priorização de ações
        \item Sensibilidade a pequenas variações no enunciado
    \end{itemize}
    \item Implicação:
    \begin{itemize}
        \item Desafio em tarefas que exigem julgamento situacional
    \end{itemize}
\end{itemize}





\section{Considerações Finais do Capítulo}
\label{subsec:fund_final}

\textbf{Ideias a abordar:}
\begin{itemize}
    \item Síntese dos conceitos:
    \begin{itemize}
        \item PMP como avaliação baseada em decisão e contexto
        \item LLMs como modelos probabilísticos com capacidades e limitações
        \item Prompt engineering como fator crítico de desempenho
    \end{itemize}
    \item Insight central:
    \begin{itemize}
        \item Há um desalinhamento potencial entre a natureza do PMP e as limitações dos LLMs
    \end{itemize}
    \item Justificativa implícita para o estudo:
    \begin{itemize}
        \item Necessidade de avaliar empiricamente esse comportamento
    \end{itemize}
    \item Transição:
    \begin{itemize}
        \item Indicar que o capítulo seguinte apresenta os trabalhos relacionados
    \end{itemize}
\end{itemize}